% Generated draft: requires mathematical and source review.
\documentclass[11pt]{article}
\usepackage[margin=25mm]{geometry}
\usepackage{fontspec}
\setmainfont{Latin Modern Roman}
\usepackage{amsmath,amssymb,mathtools,unicode-math}
\setmathfont{Latin Modern Math}
\usepackage{enumitem,needspace,xurl,hyperref}
\setlength{\parindent}{0pt}
\setlength{\parskip}{5pt}
\setlength{\emergencystretch}{3em}
\newcommand{\N}{\mathbb{N}}
\newcommand{\Z}{\mathbb{Z}}
\newcommand{\Q}{\mathbb{Q}}
\newcommand{\R}{\mathbb{R}}
\newcommand{\C}{\mathbb{C}}
\newcommand{\F}{\mathbb{F}}
\newcommand{\A}{\mathbb{A}}
\newcommand{\PP}{\mathbb{P}}
\newcommand{\E}{\mathbb{E}}
\newcommand{\eps}{\varepsilon}
\newcommand{\dd}{\,\mathrm{d}}
\begin{document}
\section*{UnsolvedMath ID: 6800017}
\textbf{Source problem code:} catalogue entry
\subsection{Definitions and mathematical statement}
In the standard lens-space setting a congruence lattice has the form
\[
L(q;s)=\{a\in\Z^d:\textstyle\sum_{j=1}^d a_js_j\equiv0\pmod q\},
\]
with \(q\ge1\) and \(\gcd(q,s_j)=1\) for the associated free cyclic sphere action. Define
\[
N_L(k)=|\{a\in L:\|a\|_1=k\}|,\qquad
\|a\|_1=\sum_j|a_j|.
\]
Norm-one isospectrality ordinarily means equality of these counts for every \(k\ge0\). A customary starred refinement counts zero coordinates as well:
\[
N_L(k,r)=|\{a\in L:\|a\|_1=k,\ |\{j:a_j=0\}|=r\}|.
\]
Equality for every \(k,r\) then implies norm-one equality. The source does not define its \(\mathrm{norm}_1^*\) notation, so this refinement is an explicit convention to verify, not an asserted recovered definition.

The corresponding round lens space is \(L(q;s)=S^{2d-1}(1)/C_q\), where a generator acts by
\[
(z_1,\ldots,z_d)\longmapsto
(e^{2\pi\mathrm i s_1/q}z_1,\ldots,e^{2\pi\mathrm i s_d/q}z_d).
\]
A pair is \(p\)-isospectral when the Hodge Laplacians \(\Delta_p=\mathrm d\delta+\delta\mathrm d\) on \(p\)-forms have identical spectra including multiplicities. All-p isospectral means this for every \(0\le p\le2d-1\).

The four requested tasks are:
\begin{enumerate}
\item Construct norm-one and starred-norm isospectral congruence lattices in all intended dimensions, with the intended nontriviality/equivalence restrictions.
\item Determine whether an all-p-isospectral family of lens spaces can contain more than two distinct examples, presumably pairwise nonisometric; a family of at least three is the literal target, not arbitrarily large families in every dimension.
\item Give a dimension-raising procedure taking all-p-isospectral lens spaces of the stated dimension \(m\) to dimension \(m+1\).
\item Establish connections with toric geometry.
\end{enumerate}
The standard real dimensions of the displayed lens spaces are odd, so task (3) cannot be read literally as ordinary real dimension plus one without an additional convention. A lattice/complex parameter increase would instead raise that real dimension by two, but this repair is not silently made. Similarly all dimensions needs a lower threshold or exceptional cases if genuinely different lattice examples are required; in dimension one the usual lattice is simply \(q\Z\).

The supplied research summary reports many constructions and links with congruence lattices but does not verify all four tasks as settled. No theorem equating the starred counts and full Hodge spectra is added without its missing hypotheses.
\subsection{Short English statement}
Construct the requested isospectral congruence lattices, seek all-form-isospectral lens-space families of size at least three, find the intended dimension-raising operation, and connect the theory to toric geometry. The starred norm and dimension conventions need clarification.
\subsection{Sources}
\begin{itemize}
\item Local UnsolvedMath registry, permanent ID \(6800017\), exact source code \texttt{catalogue entry}, ``On the Hodge spectra of lens spaces.'' Proposer Roberto Miatello, Córdoba, Argentina, 2018.
\item Morgan and Pansu, \emph{A List of Open Problems in Differential Geometry} (2018), Question environment 17; \url{https://www.imo.universite-paris-saclay.fr/~pansu/problems_MTDG.tex}. Source-TeX extraction; recorded access 2026-07-29.
\item Two source references appear only as empty ``see'' links and are missing. The supplied PARTIAL-PROGRESS summary mentions Miatello--Rossetti and unspecified coauthors, with phrases ``Lens spaces and isospectrality'' and ``p-isospectrality of lens spaces,'' but no complete publication details or exact hypotheses.
\end{itemize}
\end{document}
